\chapter{Real-полуслед и происхождение половинного веса кривизны}
\label{ch:v7-real-half-trace-curvature-weight-gate}

\section{Проверяемая гипотеза}

Предыдущий гейт получил правильный селективный гессиан для
\begin{equation}
 \mathcal S_\beta=\mathcal S_E+\beta\mathcal S_V
 \quad\text{при}\quad \beta=\frac12,
 \label{eq:v7-half-trace-target}
\end{equation}
но происхождение половины осталось открытым. Проверим, возникает ли она из
полного Real-удвоения и одного физического полуследа.

\section{Одинаковая Hodge-нормировка двух блоков}

Рёберное действие уже записано как
\begin{equation}
 \mathcal S_E=\frac12\Tr_{\mathcal K_E}\mathfrak m_E^2,
 \label{eq:v7-half-trace-edge-hodge}
\end{equation}
с точностью до постоянного вакуумного вычитания и прямых сумм, порождающих
кратности и Casimir-веса. Физическая относительная кривизна имеет ту же
структуру:
\begin{equation}
 \mathcal S_V=\frac12\left(
 \|A^*A-A_0^*A_0\|_{\rm HS}^2+
 \|AA^*-A_0A_0^*\|_{\rm HS}^2\right)
 =\frac12\Tr_{\mathcal K_V}\mathfrak m_V^2.
 \label{eq:v7-half-trace-vertex-hodge}
\end{equation}
Здесь две нормы являются двумя концами одного прямоугольного Hodge-комплекса,
а не двумя независимо взвешиваемыми действиями.

Если Real-завершение удваивает оба носителя, а один физический полуслед делит
полный след пополам, то для $X=E,V$
\begin{equation}
 \frac12\Tr_{\mathcal K_X\oplus J\mathcal K_XJ^{-1}}
 \mathfrak m_{X,\mathbb R}^2
 =\Tr_{\mathcal K_X}\mathfrak m_X^2.
 \label{eq:v7-half-trace-uniform-cancellation}
\end{equation}
Одна и та же операция умножает оба сектора на один коэффициент. Поэтому
\begin{equation}
 \boxed{\beta_{\rm Real\ half\ trace}=1,}
 \label{eq:v7-half-trace-derived-beta-one}
\end{equation}
а не $1/2$.

\section{Почему степень формы и след Клиффорда пока не помогают}

Оба текущих момента $\mathfrak m_E$ и $\mathfrak m_V$ являются конечными
эндоморфизмами степени ноль. Общая Hodge-метрика на нуль-формах и общий
пространственный след тождественного клиффордова блока действуют на них
одинаково. Множитель $1/2!$ мог бы различить нуль- и двухформенную компоненты,
но в текущем носителе ни один из двух блоков не выведен как отдельная
двухформа. Такое назначение было бы новой структурой, а не следствием
Real-удвоения.

Литература по спектральному действию подтверждает, что относительные
коэффициенты могут быть выведены общим коэффициентом теплового ядра и
клиффордовым следом после фиксации представленного исчисления, но не одним
удалением фермионного удвоения
\cite{v7-half-trace-spectral-review,v7-half-trace-wick-doubling}.

\section{Гессианный итог}

Машинный аудит возвращает для выведенного равного веса
\begin{equation}
 (n_-,n_0,n_+)_{0,\,\beta=1}=(21,0,6),
 \label{eq:v7-half-trace-beta-one-origin}
\end{equation}
тогда как условный половинный вес сохраняет
\begin{equation}
 (n_-,n_0,n_+)_{0,\,\beta=1/2}=(7,0,20).
 \label{eq:v7-half-trace-beta-half-origin}
\end{equation}
Оба веса дают положительный целевой вакуум, но только второй подавляет все
двадцать тяжёлых конкурентов уже в нуле.

\begin{theorem}[No-go Real-полуследа для относительной половины]
Полное Real-удвоение обоих Hodge-блоков и один общий физический полуслед
сохраняют их исходное отношение $1:1$. Они не выводят $\beta=1/2$.
Дополнительный полуслед только физического блока формально создаёт нужное
отношение, но эквивалентен секторному центральному весу и потому запрещён
входным правилом Тома VII.
\end{theorem}

\begin{nogocriterion}[Запрет асимметричного полуследа]
Нельзя применять полуслед повторно только к $\mathcal S_V$ или объявлять
его двухформой после просмотра гессиана. Половинный вес считается выведенным
только если общий представленный дифференциальный комплекс заранее помещает
два блока в разные степени или даёт различную клиффордову кратность.
\end{nogocriterion}

\begin{protocol}
\begin{enumerate}
 \item Real-кратность рёберного блока: \textbf{два};
 \item Real-кратность физического блока: \textbf{два};
 \item один общий физический полуслед: \textbf{сокращает оба удвоения};
 \item выведенное отношение: \textbf{$\beta=1$};
 \item его нулевая сигнатура: \textbf{$(21,0,6)$};
 \item требуемое отношение: \textbf{$\beta=1/2$};
 \item асимметричный полуслед: \textbf{ручной центральный вес};
 \item различие текущих степеней формы: \textbf{нет};
 \item различие текущих клиффордовых кратностей: \textbf{нет};
 \item итог: \textbf{Real-полуследовой маршрут закрыт};
 \item следующий гейт: \textbf{представленное исчисление степеней форм и
 клиффордова следа}.
\end{enumerate}
\end{protocol}

Машинный сертификат сохранён в
\path{s2t_v7_real_half_trace_curvature_weight_gate_results.json}.

\begin{thebibliography}{9}
\bibitem{v7-half-trace-spectral-review}
A.~Devastato, M.~A.~Kurkov, F.~Lizzi,
\emph{Spectral Noncommutative Geometry, Standard Model and all that},
Int. J. Mod. Phys. A 35 (2020), 2030013; arXiv:1906.09583.
\bibitem{v7-half-trace-wick-doubling}
F.~D'Andrea, M.~A.~Kurkov, F.~Lizzi,
\emph{Wick Rotation and Fermion Doubling in Noncommutative Geometry},
Phys. Rev. D 94 (2016), 025030; arXiv:1605.03231.
\end{thebibliography}